% ---
% school: "college"
% level: "8eme"
% section: null
% chapter: "العمليات والحساب على الأعداد الكسرية"
% source: "Axiomaths"
% title: "أولوية العمليات والكسور المركّبة"
% country: "Tunisie"
% year: 2026
% professor: "Axiomaths"
% difficulty: "Moyen"
% points: null
% tags: ["fractions", "priorites-operations", "fractions-composees"]
% ---
% Exercice 8/15 de la série 8ème « العمليات والحساب على الأعداد الكسرية » (statement + \axsol + solution). Compiler avec XeLaTeX + axiomath.sty (+ axiomaths-logo.png).
\documentclass[12pt,a4paper]{article}
\usepackage[level=8b,color=15803D]{axiomath}
\axsetup{title={العمليات والحساب على الأعداد الكسرية}, corrige}

\begin{document}

\begin{corr}[title={أولوية العمليات والكسور المركّبة}, stars=3, source={Axiomaths}]
\begin{enumerate}
  \item احسب ثم اختزل :
  \[A=\dfrac{1}{2}-\dfrac{1}{3}\times\dfrac{3}{4}+\dfrac{5}{8}\qquad B=\left(\dfrac{2}{3}-\dfrac{3}{4}\right)\div\left(\dfrac{1}{6}+\dfrac{1}{4}\right)\]
  \[C=\dfrac{\dfrac{1}{2}+\dfrac{1}{3}}{\dfrac{1}{2}-\dfrac{1}{3}}\qquad D=1-\dfrac{1}{1+\dfrac{1}{2}}\]
  \item أوجد العدد $x$ الذي يحقق $B\times x=D$.
\end{enumerate}
\axsol
\begin{prop}[أولوية العمليات]
نحسب أولا ما بين القوسين، ثم الضرب والقسمة، ثم الجمع والطرح.
\end{prop}
\begin{prop}[خط الكسر]
خط الكسر يعني القسمة : \ib{$\dfrac{a}{b}=a\div b$} ، فالكسر الأعلى يُقسم على الكسر الأسفل.
\end{prop}
\begin{enumerate}[itemsep=14pt]
  \item نطبّق الأولويات في كل حالة :
  \begin{enumerate}[itemsep=12pt]
    \item الضرب أولا :
    \[\dfrac{1}{3}\times\dfrac{3}{4}=\dfrac{3}{12}=\dfrac{1}{4}\]
    \[A=\dfrac{1}{2}-\dfrac{1}{4}+\dfrac{5}{8}=\dfrac{4}{8}-\dfrac{2}{8}+\dfrac{5}{8}=\dfrac{7}{8}\]
    \item ما بين القوسين أولا :
    \[\dfrac{2}{3}-\dfrac{3}{4}=\dfrac{8}{12}-\dfrac{9}{12}=-\dfrac{1}{12}\qquad\dfrac{1}{6}+\dfrac{1}{4}=\dfrac{2}{12}+\dfrac{3}{12}=\dfrac{5}{12}\]
    \[B=-\dfrac{1}{12}\div\dfrac{5}{12}=-\dfrac{1}{12}\times\dfrac{12}{5}=-\dfrac{1}{5}\]
    \item البسط والمقام عددان كسريان :
    \[\dfrac{1}{2}+\dfrac{1}{3}=\dfrac{5}{6}\qquad\dfrac{1}{2}-\dfrac{1}{3}=\dfrac{1}{6}\]
    \[C=\dfrac{5}{6}\div\dfrac{1}{6}=\dfrac{5}{6}\times6=5\]
    \item نبدأ من الأسفل :
    \[1+\dfrac{1}{2}=\dfrac{3}{2}\]
    \[\dfrac{1}{\dfrac{3}{2}}=1\div\dfrac{3}{2}=1\times\dfrac{2}{3}=\dfrac{2}{3}\]
    \[D=1-\dfrac{2}{3}=\dfrac{1}{3}\]
  \end{enumerate}
  \item العدد $x$ هو حاصل قسمة $D$ على $B$ :
  \[x=D\div B=\dfrac{1}{3}\div\left(-\dfrac{1}{5}\right)=\dfrac{1}{3}\times(-5)=-\dfrac{5}{3}\]
  \textbf{للتأكد :} $\left(-\dfrac{1}{5}\right)\times\left(-\dfrac{5}{3}\right)=\dfrac{1}{3}$ $\checkmark$
\end{enumerate}
\textbf{انتبه :} في $D$ ، الكسر $\dfrac{1}{1+\dfrac12}$ لا يساوي $\dfrac{1}{1}+\dfrac{1}{\dfrac12}$.
\end{corr}

\end{document}
